\section{Discussion}

The main result is an interface result.  Predicting tokenizer IDs turns the original
embedding-imitation idea into an explicit communication channel between the visual
encoder and the LLM.  It also separates two sources of success: \system{} must first
place the correct IDs on that channel; only then can \Gemma{} contribute language
knowledge.  Fluency from the downstream model cannot conceal an incorrect visual
sequence because upstream exact and edit metrics are reported independently.

The current evidence does not yet support the strongest product story---that
\Gemma{} repairs imperfect sign transcription.  It supports the prerequisite:
sign streams from held-out signers can be mapped directly into ordered, valid
\Gemma{} IDs with
substantially higher sequence accuracy than the initial decoder.  The negative
corruption audit explains why correction is absent here and produces a concrete
data requirement for the next stage: natural multiword sentences, multiple
plausible contexts, and enough redundancy that a corrupted token is inferable.

Discrete IDs trade flexibility for auditability.  Unlike continuous prompts, they
do not communicate a calibrated posterior over alternative tokens.  Future work
should compare top-$k$ or posterior-weighted interfaces, as explored in speech
prompting \citep{deng2025wav2prompt,ma2025legoslm}, while retaining a separately
scored visual transcript.  Such uncertainty may give the LLM useful alternatives
without returning to unconstrained embedding imitation.
